\section{Results and Discussion}
\label{sec:results}

\subsection{mdHODMDit Applied to the CORIA Dataset} 
\label{sec:results_coria}
This section presents the application of the multidimensional iterative Higher-Order Dynamic Mode Decomposition (mdHODMDit) algorithm to the three NH$_3$/H$_2$ DNS datasets from CORIA described in Section~\ref{sec:DATA_CORIA}. The work is carried out in collaboration with CORIA and builds upon previous developments related to the mdHODMDit methodology, particularly the work of Pablo López Salazar. Within the present study, the computational workflow was applied to the CORIA datasets, leading to the analysis of the results discussed in this section.

\subsubsection{Data Structure and Preprocessing}
\label{sec:coria_preprocessing}

The DNS output for each temperature case is stored as a four-dimensional tensor $\mathcal{T} \in \mathbb{R}^{N_{\mathrm{var}} \times N_x \times N_y \times N_t}$, where the axes correspond to the reactive scalar variables, the streamwise spatial coordinate $x$, the cross-stream coordinate $y$, and time $t$, respectively.

Before analysis, the tensor is restricted to the region of physical interest. In the streamwise direction, indices $x_{60:1000}$ are retained, corresponding to the range $[0.3, 5.0]$~cm downstream of the splitter plate, which encompasses the near-field flame stabilisation region and the initial development of the shear layer. In time, the first 74 snapshots are discarded to remove the initial transient. The temporal resolution is $\Delta t = 40~\mu$s.

For this kind of data, scaling is important due to the differences in the order of magnitude of the variables considered and the differences between their units of measurement \cite{corrochano2023higher, parente2013principal}. Since there is no universally accepted quantitative criterion to determine which scaling method is most appropriate, the analysis was performed using both auto scaling $x_{i}^{\text{scaled}} = \frac{x_i - \mu_x}{\sigma_x}$ and range scaling $x_{i}^{\text{scaled}} = \frac{x_i - x_{\min}}{x_{\max} - x_{\min}}$. The analysis yielded similar results for both scaling methods, identifying the same main frequencies.

\subsubsection{Parameter Calibration}
\label{sec:coria_params}

The analysis is performed over a grid of Hankel delay parameters
\begin{equation}
    d \in \{10,\, 15,\, 20,\, 25,\, 30,\, 35,\, 40,\, 45\}
    \label{eq:d_values}
\end{equation}
and tolerance values
\begin{equation}
    \varepsilon_{\mathrm{tol}} \in
    \{5\times10^{-3},\, 2\times10^{-2},\, 3\times10^{-2},\, 4\times10^{-2},\,
      5\times10^{-2},\, 7\times10^{-2},\, 1\times10^{-1},\, 2\times10^{-1}\}.
    \label{eq:tol_values}
\end{equation}
The same tolerance $\varepsilon_{\mathrm{tol}}$ governs both the HOSVD truncation (equation~\eqref{eq:hosvd_trunc}) and the amplitude threshold applied to retain DMD modes, so that a single parameter controls the compression level throughout the algorithm.

% -------------------------------------------------------
\subsubsection{Identification of Physical Frequencies}
\label{sec:coria_frequencies}
% -------------------------------------------------------
A DMD mode is identified as physical when its spectral peak appears across different $d$ values and it's resilient to the tolerance threshold $\varepsilon_{\mathrm{tol}}$. This is due to the fact that the coherent structures exist indipendently on how the observation matrix is built, thus should be found indipendently on how the $d$ parameter varies, while keeping in mind the previous discussion. The relative resilience of the peaks with respect to $\varepsilon_{\mathrm{tol}}$ is because relevant structures should be associated to relevant eigenvectors and not ones close to noise threshold. Frequencies are expressed in units of Strouhal number: 
\begin{equation}
    St = f \cdot \frac{\ell}{U_b},
\label{eq:strouhal}
\end{equation}
where $\ell = 3$~mm is the reference lenght of the mixing layer and  $U_b$ is the bulk speed, respectively $U_b = 21{,}34$, $24{,}29$ e $27{,}14$~m~s$^{-1}$ for cases at 1200~K, 1500~K e 1700~K.


\subsubsection{Frequency--Amplitude Spectra and associated DMD modes}
\label{sec:freq_amplitude}

The frequency--amplitude plots shown in Figures~\ref{fig:spectrum_1200K}, \ref{fig:spectrum_1500K}, and \ref{fig:spectrum_1700K} display, for each temperature case, the normalised DMD mode amplitudes $|\hat{a}_k| / \max_k |\hat{a}_k|$ as a function of Strouhal number $St$. Each marker corresponds to one $(d, \varepsilon_{\mathrm{tol}})$ pair; different colours encode different values of $d$ and different marker styles encode different tolerances. Peaks that persist across all curves can therefore be read off directly as physically robust frequencies. To showcase the coherent structures associated to these freuqencies, temperature is choosen as the rapresentative variable due to the fact that in many cases is a variable of interest for both experimental and theoretical problems. Despite this, coherent structures for every variable can be plotted.

\paragraph{Case 1200~K}
\label{sec:spectrum_1200K}

The spectrum for the 1200~K case, shown in Figure~\ref{fig:spectrum_1200K}, is characterised by a rich distribution of modes spanning a broad Strouhal range up to $St \approx 1.7$. The dominant peak is located at $St \approx 0.42$ and is reproduced consistently across all $(d, \varepsilon_{\mathrm{tol}})$ combinations, confirming its physical robustness. Secondary peaks of significant amplitude are identified at $St \approx 0.85$, $St \approx 1.1$, $St \approx 1.3$,, $St \approx 1.7$. These peaks are approximately integer multiples of the dominant frequency and therefore interpretable as higher harmonics of the fundamental shear-layer instability. 

\FloatBarrier
\begin{figure}[h!]
    \centering
    \includegraphics[width=\textwidth]{Frequency_amplitude_plot1200.png}
    \caption{Normalised DMD mode amplitudes as a function of Strouhal number for the 1200~K dataset, for $d \in \{35,40,45\}$ and $\varepsilon_{\mathrm{tol}} \in \{5\times10^{-2},\, 7\times10^{-2}\}$.}
    \label{fig:spectrum_1200K}
\end{figure}
\FloatBarrier

The associated modes to the highlighted points in Figure~\ref{fig:spectrum_1200K} are shown in Figure~\ref{fig:modes_1200K}.

\FloatBarrier
\begin{figure}[h!]
    \centering
    \begin{subfigure}[b]{0.48\textwidth}
        \centering
        \fbox{\rule{0pt}{4cm}\rule{0.9\linewidth}{0pt}}
        \caption{[PLACEHOLDER] Mode 1}
        \label{fig:modes_1200K_1}
    \end{subfigure}
    \hfill
    \begin{subfigure}[b]{0.48\textwidth}
        \centering
        \fbox{\rule{0pt}{4cm}\rule{0.9\linewidth}{0pt}}
        \caption{[PLACEHOLDER] Mode 2}
        \label{fig:modes_1200K_2}
    \end{subfigure}
    \\[1em]
    \begin{subfigure}[b]{0.48\textwidth}
        \centering
        \fbox{\rule{0pt}{4cm}\rule{0.9\linewidth}{0pt}}
        \caption{[PLACEHOLDER] Mode 3}
        \label{fig:modes_1200K_3}
    \end{subfigure}
    \hfill
    \begin{subfigure}[b]{0.48\textwidth}
        \centering
        \fbox{\rule{0pt}{4cm}\rule{0.9\linewidth}{0pt}}
        \caption{[PLACEHOLDER] Mode 4}
        \label{fig:modes_1200K_4}
    \end{subfigure}
    \caption{[PLACEHOLDER] DMD modes associated with the dominant peaks in the 1200~K spectrum.}
    \label{fig:modes_1200K}
\end{figure}
\FloatBarrier


\paragraph{Case 1500~K}
\label{sec:spectrum_1500K}

The 1500~K spectrum, shown in Figure~\ref{fig:spectrum_1500K}, exhibits a structure closely analogous to the 1200~K case, with the dominant peak again located at $St \approx 0.42$. Secondary peaks are found at $St \approx 0.78$, $St \approx 0.95$, and $St \approx 0.1$, $St \approx 1.5$. The overall behaviour is similar to the previous case reported in Fig~\ref{fig:spectrum_1200K} with the harmonic sequence translated. 
\FloatBarrier
\begin{figure}[h!]
    \centering
    \includegraphics[width=\textwidth]{Frequency_amplitude_plot1500.png}
    \caption{Normalised DMD mode amplitudes as a function of Strouhal number for the 1500~K dataset, for $d \in \{35,40,45\}$ and $\varepsilon_{\mathrm{tol}} \in \{5\times10^{-2},\, 7\times10^{-2}\}$.}
    \label{fig:spectrum_1500K}
\end{figure}
\FloatBarrier

The associated modes to the highlighted points in Figure~\ref{fig:spectrum_1500K} are shown in Figure~\ref{fig:modes_1500K}.

\FloatBarrier
\begin{figure}[h!]
    \centering
    \begin{subfigure}[b]{0.48\textwidth}
        \centering
        \fbox{\rule{0pt}{4cm}\rule{0.9\linewidth}{0pt}}
        \caption{[PLACEHOLDER] Mode 1}
        \label{fig:modes_1500K_1}
    \end{subfigure}
    \hfill
    \begin{subfigure}[b]{0.48\textwidth}
        \centering
        \fbox{\rule{0pt}{4cm}\rule{0.9\linewidth}{0pt}}
        \caption{[PLACEHOLDER] Mode 2}
        \label{fig:modes_1500K_2}
    \end{subfigure}
    \\[1em]
    \begin{subfigure}[b]{0.48\textwidth}
        \centering
        \fbox{\rule{0pt}{4cm}\rule{0.9\linewidth}{0pt}}
        \caption{[PLACEHOLDER] Mode 3}
        \label{fig:modes_1500K_3}
    \end{subfigure}
    \hfill
    \begin{subfigure}[b]{0.48\textwidth}
        \centering
        \fbox{\rule{0pt}{4cm}\rule{0.9\linewidth}{0pt}}
        \caption{[PLACEHOLDER] Mode 4}
        \label{fig:modes_1500K_4}
    \end{subfigure}
    \caption{[PLACEHOLDER] DMD modes associated with the dominant peaks in the 1500~K spectrum.}
    \label{fig:modes_1500K}
\end{figure}
\FloatBarrier


\paragraph{Case 1700~K}
\label{sec:spectrum_1700K}

The 1700~K case, shown in Figure~\ref{fig:spectrum_1700K}, presents a different spectral structure. The dominant peak remains at $St \approx 0.40$, consistent with the lower-temperature cases, but the overall spectrum is sparser: maximum relevant Strouhal number observed is $St \approx 1.25$ and fewer secondary modes are retained.

\FloatBarrier
\begin{figure}[h!]
    \centering
    \includegraphics[width=\textwidth]{Frequency_amplitude_plot1700.png}
    \caption{Normalised DMD mode amplitudes as a function of Strouhal number for the 1700~K dataset, for $d \in \{35,40,45\}$ and $\varepsilon_{\mathrm{tol}} \in \{5\times10^{-2},\, 7\times10^{-2}\}$.}
    \label{fig:spectrum_1700K}
\end{figure}
\FloatBarrier

The associated modes to the highlighted points in Figure~\ref{fig:spectrum_1700K} are shown in Figure~\ref{fig:modes_1700K}.

\FloatBarrier
\begin{figure}[h!]
    \centering
    \begin{subfigure}[b]{0.48\textwidth}
        \centering
        \fbox{\rule{0pt}{4cm}\rule{0.9\linewidth}{0pt}}
        \caption{[PLACEHOLDER] Mode 1}
        \label{fig:modes_1700K_1}
    \end{subfigure}
    \hfill
    \begin{subfigure}[b]{0.48\textwidth}
        \centering
        \fbox{\rule{0pt}{4cm}\rule{0.9\linewidth}{0pt}}
        \caption{[PLACEHOLDER] Mode 2}
        \label{fig:modes_1700K_2}
    \end{subfigure}
    \\[1em]
    \begin{subfigure}[b]{0.48\textwidth}
        \centering
        \fbox{\rule{0pt}{4cm}\rule{0.9\linewidth}{0pt}}
        \caption{[PLACEHOLDER] Mode 3}
        \label{fig:modes_1700K_3}
    \end{subfigure}
    \hfill
    \begin{subfigure}[b]{0.48\textwidth}
        \centering
        \fbox{\rule{0pt}{4cm}\rule{0.9\linewidth}{0pt}}
        \caption{[PLACEHOLDER] Mode 4}
        \label{fig:modes_1700K_4}
    \end{subfigure}
    \caption{[PLACEHOLDER] DMD modes associated with the dominant peaks in the 1700~K spectrum.}
    \label{fig:modes_1700K}
\end{figure}
\FloatBarrier

\subsubsection*{Discussion: Spectral and Spatial Evolution Across Preheat Temperatures}

The three preheat temperature cases investigated here represent distinct combustion regimes with direct relevance to the efficiency--pollutant trade-off inherent in NH$_3$/H$_2$ combustion. The 1200K case, in practical application has lower thermal efficiency, while the high temperature one, i.e. 1700~K is overall more efficient due to a more complete combustion but at the cost of significantly enhanced NOx formation. The 1500~K case sits at the boundary between these two regimes and is therefore of particular practical interest as a candidate operating point for studying the tradoff between combustion efficiency and pollutant emissions. In particular we focused on understanding the behaviour of the flow rather than the chemical one.

At 1200~k the structure of the flow found by the algorithm is harmonic, with a dominant peak at $St \approx 0.42$ is accompanied by secondary peaks at approximately integer multiples of this frequency extending up to $St \approx 1.7$. This is the sign of a strongly periodic shear layer instability, in which the primary roll up vortex and its successive pairings organize the flow in a regular pattern.

On the other hand, at 1700~K, the spectral picture changes substantially. While the dominant frequency is preserved at $St \approx 0.40$, confirming that the primary instability mechanism remains kinematically controlled by the shear between the two inlet streams, the harmonic content is severly reduced, and the overall spectrum is sparser. This evolution is physically consistent with the effect of increased thermal loading: the elevated heat release rate and associated gas expansion at 1700~K tend to suppress vortex pairing and subharmonic modes, leading to a quasi-periodic, more turbulent flow. 

The 1500~K case contains the transition between the two cases. In this spectrum we can find dominant harmonic structures as in the 1200~K, but additionally the creation of quasi-periodic frequencies that are not in a neighborhood of the main one. This quasi-periodic spectral content can be interpreted as a signature of the competing dynamics active at this operating point: the primary shear-layer instability is still present, but thermal effects are beginning to modify the nonlinear interactions.

The coherent structures corresponding to the main physical frequencies identified provide further information about the dynamics of the phoenomenon. At 1200~K, the modes with highest amplitude are spatially extended structures concentrated in the far field of the domain as shown in the Figure~DAMETTERE, reflecting the downstream convection and roll-up of large-scale vortices that have had space and time to develop along the mixing layer. At 1700~K the dominant spatial activity is shifted towards the near field, immediately downstream of the splitter plate. This can be explained by the fact that the stronger heat release ties the energetically significant dynamics to the flame attachment as shown in Figure~DA METTERE with less coherent large-scale activity further downstream. 
The transient case, i.e. 1500~K has energetically significant DMD modes in both the far field and the near field regions as shown in the DMD modes associated to the frequencies DAMETTERE and the DMD MODES DA METTERE. In this regime the instabilities are affected by both the far field convective instability and the energy release near the inlet.

\section{Conclusions}

The algorithm proposed has identified that the preheat temperature can be seen as a modulatory parameter on the shear-layer instability. The increase in temperature suppresses the secondary harmonics and shift the dynamics of the instability from the far field to the near one. The speactral profile of the flow changes from periodic to quasi periodic. 
Further investigation of the 1500~K operating point, through parametric studies and refined modal decomposition techniques, may enable the identification of conditions that simultaneously minimise NOx emissions and maximise combustion efficiency

